% ============================================================
%  PROYECTO DE DIRECCIÓN DE RECURSOS HUMANOS I
%  Doble Grado ADE-Ingenierías — Universidad de Granada
%  Empresa analizada: CaixaBank
%
%  ENTREGA 1 — Temas 2 y 3
%  Presentación oral: 10 minutos (todos los miembros deben hablar)
%
%  Para compilar localmente:
%    pdflatex main.tex  (dos veces para índices y referencias)
%  O con latexmk (recomendado):
%    latexmk -pdf main.tex
% ============================================================

\documentclass[12pt, a4paper]{article}

% ---- CODIFICACIÓN Y LENGUA ----
\usepackage[utf8]{inputenc}
\usepackage[T1]{fontenc}
\usepackage[spanish, es-tabla]{babel}

% ---- MÁRGENES ----
\usepackage[
  top=2.5cm,
  bottom=2.5cm,
  left=3cm,
  right=2.5cm
]{geometry}

% ---- TIPOGRAFÍA ----
\usepackage{lmodern}         % Fuente Latin Modern (mejor renderizado PDF)
\usepackage{microtype}       % Mejora el espaciado entre caracteres

% ---- COLORES Y GRÁFICOS ----
\usepackage[dvipsnames, table]{xcolor}
\usepackage{graphicx}
\usepackage{float}           % Control de posición de figuras [H]

% ---- TABLAS ----
\usepackage{booktabs}        % Líneas profesionales en tablas
\usepackage{longtable}       % Tablas que cruzan páginas
\usepackage{multirow}        % Celdas con múltiples filas
\usepackage{array}
\usepackage{tabularx}        % Tablas con ancho fijo y columnas flexibles

% ---- LISTAS ----
\usepackage{enumitem}
\setlist[itemize]{itemsep=2pt, topsep=4pt}
\setlist[enumerate]{itemsep=2pt, topsep=4pt}

% ---- CABECERAS Y PIES DE PÁGINA ----
\usepackage{fancyhdr}
\pagestyle{fancy}
\fancyhf{}
\fancyhead[L]{\small\textit{Dirección de Recursos Humanos I — CaixaBank}}
\fancyhead[R]{\small\textit{UGR · ADE-Ingenierías}}
\fancyfoot[C]{\thepage}
\renewcommand{\headrulewidth}{0.4pt}

% ---- FORMATO DE SECCIONES ----
\usepackage{titlesec}
\definecolor{ugrazul}{RGB}{0, 75, 135}   % Azul corporativo UGR
\titleformat{\section}{\large\bfseries\color{ugrazul}}{}{0em}{}[\titlerule]
\titleformat{\subsection}{\normalsize\bfseries\color{ugrazul!80}}{}{0em}{}
\titleformat{\subsubsection}{\normalsize\itshape}{}{0em}{}

% ---- HIPERVÍNCULOS ----
\usepackage[
  colorlinks=true,
  linkcolor=ugrazul,
  urlcolor=ugrazul,
  citecolor=ugrazul,
  pdftitle={Análisis de RRHH — CaixaBank},
  pdfauthor={Grupo RRHH UGR},
  pdfsubject={Dirección de Recursos Humanos I — Entrega 1}
]{hyperref}

% ---- TABLA DE CONTENIDOS ----
\usepackage{tocloft}
\setlength{\cftsecindent}{0pt}
\setlength{\cftsubsecindent}{1.5em}

% ---- ESPACIO ENTRE PÁRRAFOS ----
\setlength{\parskip}{6pt}
\setlength{\parindent}{0pt}

% ---- SÍMBOLOS ----
\usepackage{eurosym}             % Símbolo del euro: \euro{}

% ---- CAJAS DE NOTAS INTERNAS (para comentarios del equipo) ----
\usepackage{tcolorbox}
\tcbuselibrary{skins, breakable}

% Caja para recordatorios al equipo (se puede eliminar antes de entregar)
\newtcolorbox{notaequipo}[1][]{
  breakable,
  colback=yellow!15,
  colframe=orange!60,
  fonttitle=\bfseries\small,
  title={Nota para el equipo},
  #1
}

% ============================================================
%  DOCUMENTO
% ============================================================
\begin{document}

% ---- PORTADA ----
% ============================================================
%  PORTADA
%  Rellenad los nombres de los integrantes del grupo y el
%  nombre del/la profesor/a antes de entregar.
% ============================================================

\begin{titlepage}
  \centering

  % Logo UGR — coloca el archivo ugr_logo.png en la carpeta /imagenes/
  % Si no tienes el logo, comenta las dos líneas siguientes.
  
\includegraphics[width=0.35\textwidth]{imagenes/ugr_logo}\\[1.5cm]

  \vspace*{1.5cm}

  {\Large \textbf{Universidad de Granada}}\\[0.3cm]
  {\large Facultad de Ciencias Económicas y Empresariales}\\[0.2cm]
  {\large Doble Grado en Administración y Dirección de Empresas e Ingenierías}\\

  \vspace{1cm}
  \noindent\rule{\linewidth}{0.8pt}\\[0.5cm]

  {\LARGE \textbf{Dirección de Recursos Humanos I}}\\[0.4cm]
  {\Large Análisis de las prácticas de RRHH en una empresa real}\\[0.2cm]

  \noindent\rule{\linewidth}{0.8pt}

  \vspace{0.8cm}

  % ---- EMPRESA ----
  {\large \textbf{Empresa analizada:}}\\[0.3cm]
  {\Huge \textbf{CaixaBank}}\\[0.2cm]
  {\normalsize Sector: Servicios financieros y bancarios}\\[0.1cm]
  {\normalsize Sede central: Valencia, España}\\

  \vspace{1cm}

  % ---- ENTREGA ----
  \colorbox{ugrazul!15}{%
    \parbox{0.85\textwidth}{%
      \centering\vspace{0.3cm}
      {\large \textbf{Entrega 1 de 4} — Temas 2 y 3}\\[0.2cm]
      {\normalsize Análisis de Puestos · Planificación de RRHH}\\
      \vspace{0.3cm}
    }%
  }

  \vspace{1.5cm}

  % ---- INTEGRANTES ----
  {\large \textbf{Integrantes del grupo:}}\\[0.4cm]
  \begin{tabular}{ll}
    \toprule
    \textbf{Nombre} & \textbf{DNI} \\
    \midrule
    NOMBRE APELLIDO 1 & 00000000A \\
    NOMBRE APELLIDO 2 & 00000000B \\
    NOMBRE APELLIDO 3 & 00000000C \\
    NOMBRE APELLIDO 4 & 00000000D \\
    NOMBRE APELLIDO 5 & 00000000E \\ % Eliminar si sois 4
    \bottomrule
  \end{tabular}

  \vspace{1cm}

  % ---- DATOS ACADÉMICOS ----
  \begin{tabular}{rl}
    \textbf{Profesor/a:}   & [Nombre del/la profesor/a] \\
    \textbf{Curso:}        & 2025--2026 \\
    \textbf{Fecha:}        & \today \\
    \textbf{Convocatoria:} & Ordinaria \\
  \end{tabular}

  \vfill
\end{titlepage}

\newpage

% ---- ÍNDICE ----
\tableofcontents
\newpage

% ============================================================
%  INTRODUCCIÓN — EMPRESA Y ORGANIGRAMA
% ============================================================
% ============================================================
%  SECCIÓN 1: INTRODUCCIÓN DE LA EMPRESA Y ORGANIGRAMA
%
%  RECORDATORIO AL EQUIPO:
%  - Descripción breve pero completa: sector, tamaño, historia reciente,
%    presencia geográfica, número de empleados, cifras clave.
%  - El organigrama debe ser claro: incluir niveles jerárquicos, áreas
%    funcionales y, si es posible, el departamento de RRHH destacado.
%  - Captad cómo está estructurada la empresa formalmente vs. cómo
%    funciona realmente (cultura, informalidad, matricial vs. funcional).
%  - Si tenéis el organigrama como imagen, ponedla en /imagenes/
%    y usad \includegraphics. Si lo hacéis en LaTeX, usad tikz o forest.
% ============================================================

\section{Introducción de la empresa y organigrama}

% ---- 1.1 Descripción general ----
\subsection{Descripción general de CaixaBank}

% TODO: Ampliar con datos obtenidos de la entrevista y fuentes oficiales.
% Fuentes útiles: Memoria Anual CaixaBank, CNMV, web corporativa.

CaixaBank, S.A. es una entidad bancaria española con sede central en Valencia.
Surgida de la reorganización del Grupo ``la Caixa'', cotiza en el IBEX~35 y es
actualmente uno de los mayores grupos financieros de España y Portugal tras la
fusión con Bankia en 2021.

\begin{table}[H]
  \centering
  \caption{Datos corporativos básicos de CaixaBank (datos orientativos; actualizad con los reales)}
  \begin{tabularx}{\textwidth}{lX}
    \toprule
    \textbf{Parámetro}           & \textbf{Valor} \\
    \midrule
    Sector                       & Servicios financieros y bancarios \\
    Año de fundación             & 1904 (origen en Caja de Pensiones para la Vejez y de Ahorros) \\
    Sede central                 & Valencia, España \\
    Número de empleados          & $\sim$45.000 (grupo consolidado, 2024) \\
    Número de oficinas           & $\sim$4.000 en España \\
    Presencia internacional      & Portugal (BPI), Luxemburgo, Reino Unido, entre otros \\
    Facturación / Margen bruto   & [Dato real a incluir] \\
    \bottomrule
  \end{tabularx}
\end{table}

% ---- 1.2 Organigrama ----
\subsection{Organigrama}

% OPCIÓN A — Imagen externa (recomendada si tenéis el organigrama oficial):
% Guardad el archivo como "organigrama.png" (o .pdf) en /imagenes/ y
% descomentad las líneas siguientes:
%
% \begin{figure}[H]
%   \centering
%   \includegraphics[width=0.95\textwidth]{imagenes/organigrama}
%   \caption{Organigrama de CaixaBank (fuente: [indicar fuente])}
%   \label{fig:organigrama}
% \end{figure}

% OPCIÓN B — Descripción textual (mientras no tengáis la imagen):

A continuación se describe la estructura organizativa de CaixaBank de manera
resumida, pendiente de completar con el organigrama gráfico oficial:

\begin{itemize}
  \item \textbf{Consejo de Administración}: máximo órgano de gobierno corporativo.
  \item \textbf{Comisión Ejecutiva}: presidida por el Presidente Ejecutivo; supervisa
        la gestión diaria.
  \item \textbf{Dirección General}: organizada en grandes áreas funcionales:
    \begin{itemize}
      \item Área Financiera
      \item Área Comercial y Marketing
      \item Área de Tecnología y Datos
      \item \textbf{Área de Personas y Organización (RRHH)}
      \item Área de Riesgos
      \item Área Jurídica y Cumplimiento Normativo
    \end{itemize}
  \item \textbf{Red Territorial}: direcciones territoriales que agrupan las
        oficinas por zonas geográficas.
\end{itemize}

\begin{notaequipo}
  Completad este apartado con el organigrama real obtenido en la entrevista o
  en la web/memoria corporativa. Indicad si la estructura es \textbf{funcional},
  \textbf{divisional} o \textbf{matricial}, y si ha cambiado tras la fusión con
  Bankia. Preguntad también si existe un departamento de RRHH centralizado o
  si está descentralizado por territorios.
\end{notaequipo}

\newpage

% ============================================================
%  TEMA 2: ANÁLISIS Y DISEÑO DE PUESTOS
% ============================================================
% ============================================================
%  TEMA 2 — APARTADO 1: ANÁLISIS DE LOS PUESTOS DE TRABAJO
%
%  RECORDATORIO AL EQUIPO:
%  Las etapas del análisis de puestos (ADP) son generalmente:
%    1. Planificación del ADP (objetivos, recursos)
%    2. Preparación (diseño de instrumentos: cuestionarios, entrevistas...)
%    3. Recogida de información (observación, entrevista, cuestionario, diario...)
%    4. Análisis y tratamiento de datos
%    5. Elaboración de la descripción y especificación del puesto
%    6. Actualización y mantenimiento
%
%  PARA CADA ETAPA INDICAD:
%    - ¿La realiza CaixaBank? ¿Cómo exactamente?
%    - Ventajas e inconvenientes de su enfoque
%    - ¿Está la empresa satisfecha con ese proceso?
%    - ¿Habría maneras más eficientes de hacerlo?
% ============================================================

\section{Tema 2: Análisis de los puestos de trabajo}

\subsection{Etapas del análisis de puestos desarrolladas en CaixaBank}

El análisis de puestos de trabajo (APT) es el proceso sistemático de recopilar,
analizar y organizar información sobre los puestos de una organización con el
objetivo de elaborar descripciones y especificaciones de puestos que sirvan de
base para múltiples decisiones de RRHH (selección, formación, retribución,
evaluación del desempeño, etc.).

A continuación se detalla en qué medida CaixaBank desarrolla cada una de las
etapas clásicas del APT:

% ---- TABLA RESUMEN DE ETAPAS ----
\begin{table}[H]
  \centering
  \small
  \caption{Resumen de etapas del APT en CaixaBank}
  \begin{tabularx}{\textwidth}{p{0.28\textwidth} p{0.1\textwidth} X}
    \toprule
    \textbf{Etapa} & \textbf{¿Se realiza?} & \textbf{Observaciones} \\
    \midrule
    1. Planificación del APT    & [Sí/No/Parcial] & [Completar con información de la entrevista] \\
    2. Preparación de instrumentos & [Sí/No/Parcial] & \\
    3. Recogida de información  & [Sí/No/Parcial] & \\
    4. Análisis y tratamiento   & [Sí/No/Parcial] & \\
    5. Descripción y especificación & [Sí/No/Parcial] & \\
    6. Actualización y mantenimiento & [Sí/No/Parcial] & \\
    \bottomrule
  \end{tabularx}
\end{table}

% ---- ETAPA 1 ----
\subsubsection{Etapa 1 — Planificación del análisis de puestos}

% TODO: Describir si CaixaBank planifica formalmente el APT.
% ¿Tiene un departamento o equipo dedicado? ¿Con qué periodicidad lo actualiza?
% ¿Lo hace internamente o externaliza (consultoría)?

[Descripción a completar tras la entrevista]

\textbf{Ventajas del enfoque adoptado:}
\begin{itemize}
  \item [Completar]
\end{itemize}

\textbf{Inconvenientes / áreas de mejora:}
\begin{itemize}
  \item [Completar]
\end{itemize}

% ---- ETAPA 2 ----
\subsubsection{Etapa 2 — Preparación de los instrumentos de recogida}

% TODO: ¿Qué métodos de recogida de información utiliza CaixaBank?
% Métodos habituales: cuestionario estructurado, entrevista al ocupante,
% entrevista al supervisor, observación directa, diario de actividades,
% incidentes críticos, conferencia técnica de expertos.
% Indicad cuáles SÍ usa, cuáles NO y por qué (coste, tipo de puesto, etc.)

[Descripción a completar]

% ---- ETAPA 3 ----
\subsubsection{Etapa 3 — Recogida de información}

% TODO: ¿Quién recoge la información? ¿El propio empleado, su supervisor,
% un analista de RRHH? ¿Hay sesgo? ¿Cómo garantizan la fiabilidad?

[Descripción a completar]

% ---- ETAPAS 4, 5 y 6 ----
\subsubsection{Etapas 4, 5 y 6 — Análisis, descripción y actualización}

% TODO: ¿Existe un catálogo formal de puestos? ¿Con qué frecuencia se revisa?
% En un banco como CaixaBank la aparición de nuevas figuras (analista de datos,
% ciberseguridad, UX bancario) exige actualizaciones frecuentes.
% ¿Lo hace CaixaBank? ¿Cómo gestionan la obsolescencia de los perfiles?

[Descripción a completar]

\begin{notaequipo}
  Preguntad explícitamente en la entrevista:
  \begin{enumerate}
    \item ¿Tienen un manual o catálogo de puestos formalizado?
    \item ¿Quién es responsable de actualizar las descripciones de puesto?
    \item ¿Con qué frecuencia se revisan? ¿Se revisaron tras la fusión con Bankia?
    \item ¿Cómo recogen la información del puesto (cuestionarios, observación...)?
    \item ¿Están satisfechos con el proceso actual? ¿Qué mejorarían?
  \end{enumerate}
\end{notaequipo}

\newpage

% ============================================================
%  TEMA 2 — APARTADO 2: DESCRIPCIÓN Y ESPECIFICACIÓN DE PUESTOS
%
%  RECORDATORIO AL EQUIPO:
%  Hay que presentar TRES puestos distintos de la empresa.
%  Para cada puesto debéis elaborar:
%
%  A) DESCRIPCIÓN DEL PUESTO (el QUÉ se hace):
%     - Identificación: denominación, departamento, nivel jerárquico, superior
%       inmediato, subordinados, fecha de elaboración.
%     - Resumen del puesto (misión en 2-3 líneas).
%     - Funciones y tareas principales (ordenadas por importancia o frecuencia).
%     - Responsabilidades: personas, equipos, presupuesto, información.
%     - Condiciones de trabajo: horario, lugar, movilidad, riesgos laborales.
%     - Relaciones internas y externas.
%
%  B) ESPECIFICACIÓN DEL PUESTO (el QUIÉN lo puede hacer):
%     - Formación mínima requerida.
%     - Experiencia mínima.
%     - Conocimientos específicos.
%     - Competencias / habilidades (técnicas y conductuales).
%     - Rasgos de personalidad relevantes.
%
%  PARA CADA PUESTO VALORAD:
%  - ¿La descripción está actualizada? ¿Refleja la realidad del puesto?
%  - ¿La especificación es adecuada (ni sobre- ni infracualificación)?
%  - ¿Está la empresa satisfecha con estas fichas?
%  - ¿Podría mejorarse? ¿Cómo?
% ============================================================

\section{Descripción y especificación de tres puestos de CaixaBank}

% Elegid tres puestos de diferentes niveles jerárquicos o áreas funcionales
% para dar variedad al análisis. Ejemplos posibles:
%   - Director/a de Oficina (nivel directivo / mando intermedio)
%   - Gestor/a de Banca Privada (nivel técnico especializado)
%   - Asesor/a Comercial de Oficina (nivel operativo)
% Ajustadlos a los puestos sobre los que hayáis conseguido información real.

% ===========================================================
\subsection{Puesto 1: [Nombre del puesto — p. ej. Director/a de Oficina]}
% ===========================================================

\subsubsection{Descripción del puesto}

\begin{table}[H]
  \centering
  \small
  \caption{Ficha de identificación — Puesto 1}
  \begin{tabularx}{\textwidth}{lX}
    \toprule
    \textbf{Campo}               & \textbf{Contenido} \\
    \midrule
    Denominación del puesto      & [Director/a de Oficina] \\
    Código / referencia interna  & [Si existe] \\
    Departamento / Área          & [Red Territorial / Oficinas] \\
    Nivel jerárquico             & [Mando intermedio] \\
    Superior inmediato           & [Director/a Territorial de Zona] \\
    Número de subordinados       & [$\sim$X empleados] \\
    Fecha de elaboración         & [Mes y año de la ficha original] \\
    Última revisión              & [Fecha] \\
    \bottomrule
  \end{tabularx}
\end{table}

\textbf{Misión del puesto:}

[Resumen en 2-3 líneas sobre la razón de ser del puesto en la empresa.
Ejemplo: ``Gestionar, coordinar y liderar el equipo de la oficina bancaria,
asegurando el cumplimiento de los objetivos comerciales y la satisfacción
del cliente, en el marco de las políticas y procedimientos del Grupo CaixaBank.'']

\textbf{Funciones y tareas principales:}
\begin{enumerate}
  \item [Tarea 1 — descripción breve]
  \item [Tarea 2]
  \item [Tarea 3]
  \item [...]
\end{enumerate}

\textbf{Responsabilidades:}
\begin{itemize}
  \item \textit{Sobre personas:} [supervisión, evaluación, formación del equipo]
  \item \textit{Sobre equipos / materiales:} [...]
  \item \textit{Sobre información confidencial:} [datos de clientes, riesgo crediticio]
  \item \textit{Sobre presupuesto:} [gestión de la P\&L de la oficina]
\end{itemize}

\textbf{Condiciones de trabajo:}
\begin{itemize}
  \item Horario: [jornada completa, partido / continua, ...]
  \item Lugar: [oficina física, posible teletrabajo parcial]
  \item Movilidad / desplazamientos: [sí/no; periodicidad]
  \item Riesgos laborales: [estrés, exposición prolongada a pantallas, ...]
\end{itemize}

\subsubsection{Especificación del puesto}

\begin{table}[H]
  \centering
  \small
  \caption{Requisitos del candidato — Puesto 1}
  \begin{tabularx}{\textwidth}{lX}
    \toprule
    \textbf{Requisito}       & \textbf{Descripción} \\
    \midrule
    Formación mínima         & [Ej. Grado en ADE, Economía, Derecho o similar] \\
    Formación complementaria & [Ej. Máster en Finanzas, certificaciones bancarias (EFA, EFP)] \\
    Experiencia mínima       & [Ej. 5 años en sector bancario, 2 en gestión de equipos] \\
    Conocimientos técnicos   & [Productos bancarios, normativa MiFID II, gestión de riesgos, ...] \\
    Idiomas                  & [Español nativo; inglés nivel B2 valorable] \\
    Competencias conductuales & [Liderazgo, orientación al cliente, toma de decisiones, ...] \\
    Herramientas / software  & [CRM bancario interno, Office 365, ...] \\
    \bottomrule
  \end{tabularx}
\end{table}

\begin{notaequipo}
  Valorad si la especificación de este puesto es realista o si exige
  requisitos excesivos (sobre-cualificación) o insuficientes. ¿Cambia mucho
  la especificación real de lo que publica CaixaBank en sus ofertas de
  empleo (InfoJobs, LinkedIn, web corporativa)?
\end{notaequipo}

% ===========================================================
\subsection{Puesto 2: [Nombre del puesto — p. ej. Gestor/a de Banca Privada]}
% ===========================================================

\subsubsection{Descripción del puesto}

\begin{table}[H]
  \centering
  \small
  \caption{Ficha de identificación — Puesto 2}
  \begin{tabularx}{\textwidth}{lX}
    \toprule
    \textbf{Campo}               & \textbf{Contenido} \\
    \midrule
    Denominación del puesto      & [Gestor/a de Banca Privada] \\
    Departamento / Área          & [Banca Privada / CaixaBank Private Banking] \\
    Nivel jerárquico             & [Técnico especializado] \\
    Superior inmediato           & [Director/a de Banca Privada] \\
    Número de subordinados       & [Ninguno / gestión de cartera propia] \\
    Fecha de elaboración / revisión & [Fecha] \\
    \bottomrule
  \end{tabularx}
\end{table}

\textbf{Misión del puesto:}

[Descripción a completar]

\textbf{Funciones y tareas principales:}
\begin{enumerate}
  \item [...]
\end{enumerate}

\textbf{Responsabilidades:}
\begin{itemize}
  \item [...]
\end{itemize}

\textbf{Condiciones de trabajo:}
\begin{itemize}
  \item [...]
\end{itemize}

\subsubsection{Especificación del puesto}

\begin{table}[H]
  \centering
  \small
  \caption{Requisitos del candidato — Puesto 2}
  \begin{tabularx}{\textwidth}{lX}
    \toprule
    \textbf{Requisito}       & \textbf{Descripción} \\
    \midrule
    Formación mínima         & [Ej. Grado en ADE / Economía] \\
    Certificaciones          & [Ej. CFA, EFP, EFPA nivel III] \\
    Experiencia mínima       & [Ej. 3-5 años en gestión de patrimonios] \\
    Conocimientos técnicos   & [Mercados financieros, fiscalidad de inversiones, ...] \\
    Competencias conductuales & [Empatía, discreción, orientación al resultado, ...] \\
    \bottomrule
  \end{tabularx}
\end{table}

\begin{notaequipo}
  Comparad este puesto con el anterior. ¿Hay diferencias claras en el nivel
  de exigencia? ¿La empresa diferencia bien entre banca personal y banca privada?
\end{notaequipo}

% ===========================================================
\subsection{Puesto 3: [Nombre del puesto — p. ej. Asesor/a Comercial de Oficina]}
% ===========================================================

\subsubsection{Descripción del puesto}

\begin{table}[H]
  \centering
  \small
  \caption{Ficha de identificación — Puesto 3}
  \begin{tabularx}{\textwidth}{lX}
    \toprule
    \textbf{Campo}               & \textbf{Contenido} \\
    \midrule
    Denominación del puesto      & [Asesor/a Comercial de Oficina] \\
    Departamento / Área          & [Red Comercial] \\
    Nivel jerárquico             & [Operativo / técnico base] \\
    Superior inmediato           & [Director/a de Oficina] \\
    Número de subordinados       & [Ninguno] \\
    Fecha de elaboración / revisión & [Fecha] \\
    \bottomrule
  \end{tabularx}
\end{table}

\textbf{Misión del puesto:}

[Descripción a completar]

\textbf{Funciones y tareas principales:}
\begin{enumerate}
  \item [...]
\end{enumerate}

\textbf{Responsabilidades:}
\begin{itemize}
  \item [...]
\end{itemize}

\textbf{Condiciones de trabajo:}
\begin{itemize}
  \item [...]
\end{itemize}

\subsubsection{Especificación del puesto}

\begin{table}[H]
  \centering
  \small
  \caption{Requisitos del candidato — Puesto 3}
  \begin{tabularx}{\textwidth}{lX}
    \toprule
    \textbf{Requisito}       & \textbf{Descripción} \\
    \midrule
    Formación mínima         & [Ej. FP Superior en Administración, Grado ADE o similar] \\
    Certificaciones          & [EFA, EFB o formación interna CaixaBank] \\
    Experiencia mínima       & [0-2 años; puede ser puesto de entrada] \\
    Conocimientos técnicos   & [Productos bancarios básicos, normativa, CRM, ...] \\
    Competencias conductuales & [Orientación al cliente, proactividad, comunicación, ...] \\
    \bottomrule
  \end{tabularx}
\end{table}

\begin{notaequipo}
  Este suele ser el puesto de entrada más común. Analizad si la rotación
  es alta en este nivel y si eso refleja problemas en la especificación
  (se selecciona a perfiles que luego no encajan) o en las condiciones del puesto.
  ¿Cuánto tarda CaixaBank en cubrir una vacante de este tipo?
\end{notaequipo}

\newpage

% ============================================================
%  TEMA 2 — APARTADO 3: FORMAS DE REDISEÑO DE PUESTOS
%
%  RECORDATORIO AL EQUIPO:
%  Las principales formas de rediseño de puestos son:
%
%    1. ROTACIÓN DE PUESTOS (job rotation): los empleados rotan entre puestos.
%       Ventajas: polivalencia, motivación, reducción de la monotonía.
%       Inconvenientes: costes de formación, pérdida de especialización.
%
%    2. AMPLIACIÓN DE TAREAS (job enlargement / ensanchamiento horizontal):
%       se añaden tareas del mismo nivel de complejidad.
%       Ventajas: reduce la monotonía, mayor variedad.
%       Inconvenientes: puede aumentar la carga sin añadir motivación real.
%
%    3. ENRIQUECIMIENTO DE PUESTOS (job enrichment / Hackman & Oldham):
%       se añaden responsabilidades de mayor complejidad (carga vertical).
%       Las cinco características centrales del modelo:
%         - Variedad de habilidades
%         - Identidad de la tarea
%         - Significado de la tarea
%         - Autonomía
%         - Retroalimentación
%       Ventajas: mayor motivación intrínseca, satisfacción, compromiso.
%       Inconvenientes: requiere empleados que quieran / puedan asumir más.
%
%    4. EQUIPOS DE TRABAJO AUTOGESTIONADOS (self-managed teams):
%       grupos con alta autonomía para organizar su propio trabajo.
%
%    5. TELETRABAJO / TRABAJO FLEXIBLE (flexibilización del puesto):
%       modificación del lugar y/o horario de trabajo.
%
%    6. JOB CRAFTING: el empleado adapta el puesto a sus fortalezas
%       de forma proactiva (tendencia emergente).
%
%  PARA CADA FORMA INDICAD:
%  - ¿La aplica CaixaBank? ¿En qué puestos / áreas?
%  - Ventajas e inconvenientes concretos observados.
%  - ¿Está la empresa satisfecha?
%  - ¿Qué formas serían interesantes aplicar y por qué?
% ============================================================

\section{Formas de rediseño de puestos en CaixaBank}

\subsection{Formas de rediseño aplicadas actualmente}

% ---- TABLA RESUMEN ----
\begin{table}[H]
  \centering
  \small
  \caption{Aplicación de formas de rediseño de puestos en CaixaBank}
  \begin{tabularx}{\textwidth}{p{0.28\textwidth} p{0.12\textwidth} X}
    \toprule
    \textbf{Forma de rediseño}      & \textbf{¿Se aplica?} & \textbf{Descripción y observaciones} \\
    \midrule
    Rotación de puestos             & [Sí/No/Parcial] & \\
    Ampliación de tareas            & [Sí/No/Parcial] & \\
    Enriquecimiento de puestos      & [Sí/No/Parcial] & \\
    Equipos autogestionados         & [Sí/No/Parcial] & \\
    Teletrabajo / trabajo flexible  & [Sí/No/Parcial] & \\
    Job crafting                    & [Sí/No/Parcial] & \\
    \bottomrule
  \end{tabularx}
\end{table}

% ---- ROTACIÓN ----
\subsubsection{Rotación de puestos (\textit{job rotation})}

% TODO: ¿Existe un programa formal de rotación en CaixaBank?
% Es habitual en banca para desarrollar directivos y gestores con visión global.
% ¿Se rota entre oficinas, entre departamentos, internacionalmente (BPI)?

[Descripción a completar]

\textbf{Ventajas observadas en CaixaBank:}
\begin{itemize}
  \item [...]
\end{itemize}

\textbf{Inconvenientes observados:}
\begin{itemize}
  \item [...]
\end{itemize}

% ---- AMPLIACIÓN ----
\subsubsection{Ampliación de tareas (\textit{job enlargement})}

% TODO: Con la digitalización bancaria, los asesores de oficina han asumido
% tareas que antes hacían otros departamentos (apertura online, soporte digital...).
% ¿Esto genera sobrecarga? ¿Hay compensación retributiva?

[Descripción a completar]

% ---- ENRIQUECIMIENTO ----
\subsubsection{Enriquecimiento de puestos (\textit{job enrichment})}

% TODO: ¿Se da autonomía a los asesores para tomar decisiones sobre productos?
% ¿Tienen feedback claro sobre su impacto en los resultados?
% El modelo de Hackman y Oldham es muy útil para evaluar si el diseño actual
% de los puestos genera motivación intrínseca.

[Descripción a completar]

% ---- TELETRABAJO ----
\subsubsection{Teletrabajo y trabajo flexible}

% TODO: Tras la pandemia COVID-19, CaixaBank implementó políticas de teletrabajo.
% ¿Sigue vigente? ¿Para qué puestos? ¿Con qué porcentaje de jornada?
% ¿Ha mejorado la satisfacción y productividad? ¿Ha generado problemas
% de desconexión del equipo o de la cultura organizacional?

[Descripción a completar]

\subsection{Formas de rediseño recomendadas para CaixaBank}

% TODO: Proponed al menos 2-3 formas que no se apliquen o que se apliquen
% de forma mejorable. Justificad por qué serían beneficiosas.
% Ejemplos posibles:
%   - Impulsar el job crafting en puestos muy regulados para mejorar motivación.
%   - Aplicar equipos autogestionados en áreas de innovación/tecnología.
%   - Revisar el job enlargement para evitar sobrecarga sin contraprestación.

\begin{notaequipo}
  Preguntad en la entrevista:
  \begin{enumerate}
    \item ¿Existe algún programa formal de rotación? ¿Tiene criterios claros?
    \item ¿Los empleados de oficina han asumido más tareas en los últimos años?
          ¿Se ha revisado su carga de trabajo y retribución?
    \item ¿Tienen autonomía para adaptar su forma de trabajo (\textit{job crafting})?
    \item ¿Cuál es la política actual de teletrabajo? ¿Están satisfechos con ella?
    \item ¿Se mide el impacto de los rediseños en productividad y satisfacción?
  \end{enumerate}
\end{notaequipo}

\newpage

% ============================================================
%  TEMA 2 — APARTADO 4: NUEVAS TENDENCIAS EN ADP — POR COMPETENCIAS
%
%  RECORDATORIO AL EQUIPO:
%  El Análisis y Diseño de Puestos (ADP) por competencias supone un cambio
%  de paradigma respecto al ADP tradicional:
%
%  ADP TRADICIONAL:
%    - Se centra en TAREAS y RESPONSABILIDADES del puesto.
%    - El resultado es una descripción de puesto estática.
%    - El encaje persona-puesto se hace comparando requisitos formales
%      (titulación, experiencia) con el perfil del candidato.
%
%  ADP POR COMPETENCIAS:
%    - Se centra en los COMPORTAMIENTOS OBSERVABLES que llevan al éxito.
%    - Identifica competencias clave (técnicas + conductuales/transversales).
%    - Es más flexible: la misma competencia puede ser válida en varios puestos.
%    - Facilita la gestión integrada del talento (selección, formación,
%      evaluación del desempeño, planes de carrera, retribución) bajo un
%      marco de competencias común.
%    - Permite adaptar rápidamente los perfiles a entornos cambiantes (VUCA).
%
%  OTRAS TENDENCIAS:
%    - ADP en entornos ágiles (puestos definidos por roles y sprints, no tareas).
%    - Uso de People Analytics e IA para rediseñar puestos basándose en datos.
%    - Gestión por proyectos vs. gestión por puestos (modelos gig/fluid).
%    - ADP continuo: revisión periódica sistemática en lugar de actualización
%      puntual y reactiva.
%    - Wellbeing y diseño de puestos orientado al bienestar del empleado.
%
%  PARA ESTA SECCIÓN INDICAD:
%  - ¿Aplica CaixaBank el ADP por competencias? ¿Tiene un diccionario de
%    competencias corporativo?
%  - ¿Cómo integra el enfoque de competencias en selección, formación y evaluación?
%  - ¿Usa People Analytics para analizar o rediseñar puestos?
%  - ¿Está la empresa satisfecha con su modelo? ¿Qué tendencias le faltan adoptar?
% ============================================================

\section{Nuevas tendencias en análisis y diseño de puestos: ADP por competencias}

\subsection{Del enfoque tradicional al enfoque por competencias}

El análisis y diseño de puestos (ADP) ha evolucionado desde el modelo
clásico, centrado en la descripción de tareas y responsabilidades, hacia
un enfoque basado en \textbf{competencias}, que pone el foco en los
comportamientos observables y medibles que conducen a un desempeño
excelente en el puesto.

\begin{table}[H]
  \centering
  \small
  \caption{Comparativa entre el ADP tradicional y el ADP por competencias}
  \begin{tabularx}{\textwidth}{p{0.23\textwidth} X X}
    \toprule
    \textbf{Dimensión}       & \textbf{ADP Tradicional} & \textbf{ADP por Competencias} \\
    \midrule
    Unidad de análisis       & Tareas y responsabilidades & Competencias y comportamientos \\
    Resultado                & Descripción de puesto estática & Perfil de competencias dinámico \\
    Orientación              & Presente                  & Presente y futuro \\
    Flexibilidad             & Baja                      & Alta \\
    Integración con otros procesos RRHH & Limitada   & Total (marco integrado) \\
    Adaptación al cambio     & Reactiva                  & Proactiva \\
    \bottomrule
  \end{tabularx}
\end{table}

\subsection{Aplicación del ADP por competencias en CaixaBank}

% TODO: CaixaBank, como gran empresa del IBEX 35, es probable que tenga
% un diccionario de competencias corporativo. Investigad:
%   - ¿Existe un modelo de competencias formalizado?
%   - ¿Cuáles son las competencias corporativas (transversales a todos los puestos)?
%   - ¿Hay competencias específicas por nivel o por familia de puestos?
%   - ¿Cómo se usa el modelo en selección, evaluación del desempeño y formación?

\subsubsection{Modelo de competencias de CaixaBank}

[Descripción a completar con la información de la entrevista y/o la
web corporativa de CaixaBank (sección de empleo / cultura)]

Ejemplos de competencias frecuentes en entidades bancarias (ajustar a CaixaBank):
\begin{itemize}
  \item \textbf{Orientación al cliente}: escucha activa, empatía, resolución
        de problemas del cliente.
  \item \textbf{Orientación a resultados}: planificación, seguimiento de objetivos,
        toma de decisiones basada en datos.
  \item \textbf{Trabajo en equipo y colaboración}: comunicación, confianza,
        gestión de conflictos.
  \item \textbf{Adaptabilidad e innovación}: apertura al cambio, aprendizaje
        continuo, pensamiento creativo.
  \item \textbf{Integridad y cumplimiento}: ética profesional, respeto normativo,
        confidencialidad.
\end{itemize}

\subsubsection{Integración del modelo de competencias en los procesos de RRHH}

% TODO: ¿Cómo se usa el diccionario de competencias en:
%   - Selección (entrevista por competencias, assesment centres)?
%   - Evaluación del desempeño (evaluación 360º, por objetivos + competencias)?
%   - Formación y desarrollo (planes de desarrollo individuales)?
%   - Planes de carrera y sucesión?

[Descripción a completar]

\subsection{Otras tendencias emergentes en ADP}

\subsubsection{People Analytics aplicado al análisis de puestos}

% TODO: ¿Usa CaixaBank datos para analizar patrones de desempeño, rotación
% o satisfacción por tipo de puesto? ¿Tienen un equipo de People Analytics?
% En una empresa de su tamaño es probable que sí. ¿Qué sistemas usan (SAP,
% Workday, Cornerstone...)?

[Descripción a completar]

\subsubsection{ADP en entornos ágiles}

% TODO: Especialmente relevante en el área de tecnología y digitalización
% de CaixaBank. ¿Trabajan con metodologías ágiles (Scrum, SAFe)?
% ¿Cómo afecta eso al diseño de puestos (roles vs. puestos fijos)?

[Descripción a completar]

\subsubsection{Diseño de puestos y bienestar del empleado (\textit{wellbeing})}

% TODO: CaixaBank ha publicado informes de sostenibilidad con compromisos
% en bienestar del empleado. ¿Se refleja eso en el diseño de los puestos?
% ¿Tienen indicadores de bienestar vinculados al diseño de puestos?

[Descripción a completar]

\begin{notaequipo}
  Preguntad en la entrevista:
  \begin{enumerate}
    \item ¿Tienen un diccionario o modelo de competencias corporativo?
          ¿Podéis verlo o que os lo describan?
    \item ¿Cómo se usa ese modelo en la selección y en la evaluación del desempeño?
    \item ¿Tienen un equipo o función de People Analytics?
    \item ¿Trabajan con metodologías ágiles en alguna área? ¿Cómo ha cambiado
          el concepto de ``puesto'' en esas áreas?
    \item ¿Qué tendencias en ADP creen que les falta adoptar?
  \end{enumerate}
\end{notaequipo}

\newpage

% ============================================================
%  TEMA 3: PLANIFICACIÓN DE RRHH
% ============================================================
% ============================================================
%  TEMA 3 — APARTADO 1: PLANIFICACIÓN DE RRHH O DE PLANTILLA
%
%  RECORDATORIO AL EQUIPO:
%  La planificación de RRHH (o planificación de plantilla) consiste en
%  anticipar las necesidades futuras de personal de la organización y
%  definir las acciones necesarias para satisfacerlas.
%
%  NIVELES DE PLANIFICACIÓN:
%    1. OPERATIVA (corto plazo, 1 año): cubrir necesidades inmediatas,
%       absentismo, bajas, picos de demanda.
%    2. TÁCTICA (medio plazo, 1-3 años): ajuste de plantilla, formación,
%       movilidad interna.
%    3. ESTRATÉGICA (largo plazo, +3 años): alineada con la estrategia
%       de negocio; transforma la estructura de la plantilla para los
%       retos futuros (digitalización, nuevos mercados, etc.).
%
%  PLANIFICACIÓN ESTRATÉGICA DE RRHH:
%    - Está integrada con el plan estratégico de la empresa.
%    - Considera escenarios futuros (tecnología, regulación, demografía).
%    - Incluye planificación de la sucesión de posiciones clave.
%    - Se mide con indicadores (KPIs de plantilla, turnover, etc.).
%
%  PARA ESTA SECCIÓN INDICAD:
%  - ¿Existe un proceso formal de planificación de plantilla en CaixaBank?
%  - ¿Es operativa, táctica o estratégica (o las tres)?
%  - ¿Está integrada con el plan estratégico del grupo?
%  - ¿Está la empresa satisfecha? ¿Qué mejoraría?
% ============================================================

\section{Tema 3: Planificación de recursos humanos}

\subsection{Existencia y naturaleza de la planificación de RRHH en CaixaBank}

% TODO: CaixaBank publica planes estratégicos trienales (p.ej. Plan Estratégico
% 2022-2024). Comprobad si ese plan incluye objetivos específicos de gestión
% de personas / plantilla. La sección de Sostenibilidad de la Memoria Anual
% también suele incluir datos de plantilla y objetivos de RRHH.

La planificación de recursos humanos es el proceso mediante el cual la
organización prevé sus necesidades futuras de personal y define las acciones
necesarias —contratación, formación, reducción, reconversión— para alinear
la dotación de personal con los objetivos estratégicos del negocio.

\subsubsection{¿Existe planificación formal de RRHH en CaixaBank?}

[Descripción a completar con la información obtenida en la entrevista
y en los documentos públicos de CaixaBank (Memoria Anual, Plan Estratégico)]

\subsubsection{¿Es estratégica la planificación de RRHH?}

% Para que la planificación sea estratégica debe:
%   a) Estar integrada con el plan estratégico de negocio.
%   b) Tener un horizonte temporal de al menos 3 años.
%   c) Considerar escenarios de futuro (no solo extrapolar el pasado).
%   d) Involucrar a la alta dirección, no solo al departamento de RRHH.
%   e) Contar con indicadores de seguimiento (KPIs de plantilla).

\begin{table}[H]
  \centering
  \small
  \caption{Características de la planificación estratégica de RRHH en CaixaBank}
  \begin{tabularx}{\textwidth}{p{0.45\textwidth} p{0.12\textwidth} X}
    \toprule
    \textbf{Característica}                          & \textbf{¿Se cumple?} & \textbf{Observaciones} \\
    \midrule
    Integración con plan estratégico de negocio      & [Sí/No/Parcial] & \\
    Horizonte temporal $\geq$ 3 años                 & [Sí/No/Parcial] & \\
    Consideración de escenarios de futuro            & [Sí/No/Parcial] & \\
    Involucración de la alta dirección               & [Sí/No/Parcial] & \\
    KPIs de seguimiento definidos                    & [Sí/No/Parcial] & \\
    Planificación de la sucesión                     & [Sí/No/Parcial] & \\
    \bottomrule
  \end{tabularx}
\end{table}

\subsubsection{Impacto de la fusión con Bankia en la planificación de plantilla}

% La fusión CaixaBank-Bankia (2021) fue la mayor fusión bancaria de la historia
% de España. Implicó un ajuste de plantilla significativo (ERE 2021).
% Analizad cómo afectó a la planificación de RRHH: ¿fue reactiva o proactiva?
% ¿Cómo se gestionó la integración de plantillas?

[Descripción a completar]

\begin{notaequipo}
  Preguntad en la entrevista:
  \begin{enumerate}
    \item ¿Existe un plan de plantilla formal? ¿Quién lo elabora y con qué periodicidad?
    \item ¿Está vinculado al plan estratégico del grupo?
    \item ¿Cómo se gestionó la planificación de plantilla tras la fusión con Bankia?
    \item ¿Tienen planificación de sucesión para posiciones clave?
    \item ¿Qué KPIs de plantilla monitorizan (ratio empleado/oficina, rotación,
          cobertura de vacantes, etc.)?
  \end{enumerate}
\end{notaequipo}

\newpage

% ============================================================
%  TEMA 3 — APARTADO 2: ANÁLISIS DE OFERTA Y DEMANDA DE RRHH
%
%  RECORDATORIO AL EQUIPO:
%
%  ANÁLISIS DE DEMANDA (¿cuántos empleados necesitamos?):
%  Métodos cuantitativos:
%    - Análisis de series temporales (extrapolación de tendencias históricas).
%    - Ratios de productividad (unidades producidas / empleado).
%    - Análisis de regresión (demanda de empleo en función de variables del negocio).
%    - Modelos econométricos.
%  Métodos cualitativos:
%    - Técnica Delphi (consenso de expertos internos).
%    - Juicio directivo (managers estiman sus necesidades).
%    - Escenarios (qué pasa si... → planificación por escenarios).
%
%  ANÁLISIS DE OFERTA (¿cuántos empleados tenemos / tendremos?):
%  Oferta interna (inventario de la plantilla actual):
%    - Inventario de habilidades / perfiles.
%    - Matrices de Markov / tablas de transición (probabilidad de que un
%      empleado permanezca, sea promovido, se vaya...).
%    - Análisis de la rotación y del absentismo.
%    - Mapas de reemplazamiento / planes de sucesión.
%  Oferta externa (mercado laboral):
%    - Análisis del mercado de trabajo local / sectorial.
%    - Datos del INE, SEPE, convenios colectivos.
%    - Benchmarking con otras entidades financieras.
%
%  PARA ESTA SECCIÓN INDICAD:
%  - ¿Qué métodos de análisis de demanda usa CaixaBank?
%  - ¿Qué métodos de análisis de oferta usa?
%  - ¿Cuáles serían convenientes aplicar que no usa actualmente?
%  - Ventajas e inconvenientes del enfoque actual.
% ============================================================

\section{Análisis de oferta y demanda de recursos humanos}

\subsection{Análisis de la demanda de RRHH}

La previsión de la demanda de personal permite a CaixaBank anticipar cuántos
y qué tipo de empleados necesitará en el futuro para alcanzar sus objetivos
de negocio.

\subsubsection{Métodos utilizados para la previsión de la demanda}

\begin{table}[H]
  \centering
  \small
  \caption{Métodos de previsión de la demanda de RRHH en CaixaBank}
  \begin{tabularx}{\textwidth}{p{0.28\textwidth} p{0.12\textwidth} X}
    \toprule
    \textbf{Método}                  & \textbf{¿Se usa?} & \textbf{Descripción y observaciones} \\
    \midrule
    Series temporales / tendencias   & [Sí/No/Parcial] & \\
    Ratios de productividad          & [Sí/No/Parcial] & \\
    Regresión / modelos cuantitativos & [Sí/No/Parcial] & \\
    Técnica Delphi                   & [Sí/No/Parcial] & \\
    Juicio directivo                 & [Sí/No/Parcial] & \\
    Planificación por escenarios     & [Sí/No/Parcial] & \\
    \bottomrule
  \end{tabularx}
\end{table}

[Descripción detallada de los métodos más relevantes que usa CaixaBank]

\textbf{Métodos que serían convenientes aplicar:}

% TODO: Basándoos en el tamaño y complejidad de CaixaBank, ¿qué métodos
% más avanzados podrían añadir valor? Ejemplos:
%   - Modelos de regresión múltiple incorporando variables macroeconómicas
%     (volumen de crédito, tipos de interés, PIB) para prever necesidades.
%   - People Analytics con machine learning para predecir rotación.
%   - Planificación por escenarios ante disrupciones tecnológicas (IA, blockchain).

\begin{itemize}
  \item [Método 1 — justificación]
  \item [Método 2 — justificación]
\end{itemize}

\subsection{Análisis de la oferta de RRHH}

\subsubsection{Oferta interna: inventario de la plantilla}

% TODO: ¿Tiene CaixaBank un sistema de inventario de habilidades/competencias?
% ¿Usan matrices de Markov o tablas de transición para modelar flujos internos?
% ¿Tienen mapas de reemplazamiento / planes de sucesión formales?

[Descripción a completar]

\subsubsection{Oferta externa: análisis del mercado de trabajo}

% TODO: ¿Cómo monitoriza CaixaBank el mercado laboral financiero en España?
% ¿Hacen benchmarking de salarios con competidores (Santander, BBVA, Sabadell)?
% ¿Colaboran con universidades para captar talento (relaciones con la UGR,
% prácticas, campus recruitment)?

[Descripción a completar]

\begin{notaequipo}
  Preguntad en la entrevista:
  \begin{enumerate}
    \item ¿Con qué métodos cuantifican sus necesidades futuras de personal?
    \item ¿Tienen un inventario actualizado de competencias de la plantilla?
    \item ¿Usan modelos matemáticos o estadísticos para prever la oferta interna
          (p.ej. matrices de transición / Markov)?
    \item ¿Cómo monitorizan el mercado laboral externo?
    \item ¿Cuál es su tasa de rotación actual? ¿La consideran alta o baja?
    \item ¿Han tenido dificultades para cubrir ciertos perfiles en los últimos años?
          ¿Cuáles? (Probablemente perfiles tecnológicos: Data Scientists, Ingenieros,
          especialistas en ciberseguridad)
  \end{enumerate}
\end{notaequipo}

\newpage

% ============================================================
%  TEMA 3 — APARTADO 3: EQUILIBRIO ENTRE OFERTA Y DEMANDA DE RRHH
%
%  RECORDATORIO AL EQUIPO:
%  Tras comparar la oferta y la demanda de RRHH se pueden dar tres situaciones:
%
%  1. OFERTA = DEMANDA → Equilibrio. Situación ideal pero poco frecuente.
%     Acciones: mantenimiento, formación para adaptarse a cambios menores.
%
%  2. DEMANDA > OFERTA → Déficit de personal. Necesidad de incorporar.
%     Acciones: reclutamiento externo, promoción interna, horas extra,
%     trabajo temporal, externalización (outsourcing), reincorporación de
%     exempleados (boomerang employees), automatización de tareas para
%     reducir necesidades.
%
%  3. OFERTA > DEMANDA → Superávit de personal. Necesidad de reducir.
%     Acciones: ERE, ERTE, jubilaciones anticipadas, congelación de
%     contrataciones, reducción de jornada, recolocación (outplacement),
%     formación para reasignar a otras áreas. (Se desarrolla en el
%     apartado de procesos sustractivos, sección 10)
%
%  PARA ESTA SECCIÓN INDICAD:
%  - ¿Cuál es la situación de CaixaBank actualmente (déficit/equilibrio/superávit)?
%  - ¿Ha variado en los últimos años (especialmente tras la fusión)?
%  - ¿Qué acciones concretas ha tomado para lograr el equilibrio?
%  - ¿Está satisfecha con los resultados de esas acciones?
% ============================================================

\section{Equilibrio entre oferta y demanda de recursos humanos}

\subsection{Situación actual de CaixaBank}

% TODO: Tras el ERE de 2021 (fusión con Bankia, con salida de ~6.000 empleados),
% la plantilla se ha ido ajustando. Pero al mismo tiempo la digitalización
% crea demanda de nuevos perfiles tecnológicos que no siempre se pueden cubrir
% internamente. Puede haber simultáneamente superávit en ciertos perfiles
% tradicionales (cajeros, operadores de red) y déficit en perfiles digitales.

[Descripción a completar]

\begin{table}[H]
  \centering
  \small
  \caption{Situación de equilibrio oferta-demanda en CaixaBank por tipo de perfil}
  \begin{tabularx}{\textwidth}{p{0.30\textwidth} p{0.20\textwidth} X}
    \toprule
    \textbf{Familia de puestos}    & \textbf{Situación} & \textbf{Acciones adoptadas} \\
    \midrule
    Perfiles bancarios tradicionales (oficina) & [Superávit/Déficit/Equilibrio] & \\
    Perfiles tecnológicos (datos, ciberseg., digital) & [Superávit/Déficit/Equilibrio] & \\
    Perfiles de gestión y dirección  & [Superávit/Déficit/Equilibrio] & \\
    Perfiles de cumplimiento y riesgo & [Superávit/Déficit/Equilibrio] & \\
    \bottomrule
  \end{tabularx}
\end{table}

\subsection{Acciones de ajuste ante déficit de personal}

% TODO: En qué áreas hay escasez y cómo la cubre CaixaBank.
% ¿Recurre a reclutamiento externo? ¿A la promoción interna? ¿A la formación
% de recualificación (reskilling)?

\begin{itemize}
  \item \textbf{Reclutamiento externo:} [descripción]
  \item \textbf{Promoción interna:} [descripción]
  \item \textbf{Reskilling / upskilling:} [¿tienen programas de recualificación?]
  \item \textbf{Trabajo temporal o externalización:} [descripción]
  \item \textbf{Acuerdos con universidades / campus recruitment:} [descripción]
\end{itemize}

\subsection{Acciones de ajuste ante superávit de personal}

% Se desarrollará con mayor detalle en la sección de Procesos Sustractivos (sección 10).
% Aquí haced un resumen de los mecanismos de reducción utilizados.

\begin{itemize}
  \item \textbf{ERE 2021 (fusión Bankia):} [número de afectados, condiciones]
  \item \textbf{Jubilaciones anticipadas:} [descripción]
  \item \textbf{Recolocación interna (movilidad geográfica / funcional):} [descripción]
  \item \textbf{Planes de formación para reconversión profesional:} [descripción]
\end{itemize}

\begin{notaequipo}
  Preguntad en la entrevista:
  \begin{enumerate}
    \item ¿En qué perfiles sienten más escasez de candidatos en el mercado?
    \item ¿Cómo afrontaron el exceso de plantilla tras la fusión con Bankia?
    \item ¿Han implementado programas de reskilling o upskilling para
          reconvertir perfiles tradicionales en perfiles digitales?
    \item ¿Cómo gestionan las situaciones de superávit temporal
          sin recurrir a despidos (ERTE, reducción de jornada...)?
    \item ¿Está la plantilla actual equilibrada o hay áreas con tensión?
  \end{enumerate}
\end{notaequipo}

\newpage

% ============================================================
%  TEMA 3 — APARTADO 4: CRITERIOS PARA EVALUAR LA PLANIFICACIÓN
%                       Y SISTEMAS DE INFORMACIÓN DE RRHH (SIRH)
%
%  RECORDATORIO AL EQUIPO:
%
%  CRITERIOS DE EVALUACIÓN DE LA PLANIFICACIÓN DE RRHH:
%  La planificación de RRHH se evalúa mediante indicadores (KPIs) que miden
%  su eficacia y eficiencia. Los más habituales son:
%
%    Indicadores de flujo de plantilla:
%      - Tasa de rotación (turnover rate): % de bajas sobre plantilla media.
%      - Tasa de absentismo: horas no trabajadas / horas contratadas.
%      - Tiempo medio de cobertura de vacantes (time-to-fill).
%      - Tasa de promoción interna.
%
%    Indicadores de calidad del ajuste:
%      - Porcentaje de posiciones clave cubiertas por planes de sucesión.
%      - Tasa de retención en el primer año.
%      - Tasa de superación del período de prueba.
%
%    Indicadores de coste:
%      - Coste medio de contratación (cost-per-hire).
%      - Coste del absentismo.
%      - Retorno de la inversión en formación (ROI de formación).
%
%    Indicadores de previsión:
%      - Desviación entre plantilla planificada y real.
%      - % de posiciones cubiertas internamente vs. externamente.
%
%  SISTEMAS DE INFORMACIÓN DE RRHH (SIRH / HRIS):
%  Un SIRH es una plataforma de software que integra y gestiona los datos
%  de los empleados y los procesos de RRHH. Sus módulos típicos son:
%    - Administración de personal (contratos, nóminas, seguridad social).
%    - Gestión del tiempo y asistencia.
%    - Selección y reclutamiento (ATS - Applicant Tracking System).
%    - Formación y desarrollo (LMS - Learning Management System).
%    - Evaluación del desempeño.
%    - Planificación de plantilla y People Analytics.
%
%  Los sistemas más conocidos del mercado:
%    SAP SuccessFactors, Workday, Oracle HCM, Cornerstone OnDemand,
%    Meta4 / NGA HR (muy usado en España), A3 RRHH.
%
%  PARA ESTA SECCIÓN INDICAD:
%  - ¿Qué KPIs usa CaixaBank para medir la planificación de RRHH?
%  - ¿Qué sistema(s) SIRH utiliza y qué módulos tiene activos?
%  - ¿Cómo se podría mejorar el sistema o los criterios de evaluación?
% ============================================================

\section{Criterios de evaluación de la planificación y sistemas de información de RRHH}

\subsection{Criterios e indicadores para evaluar la planificación de RRHH}

La eficacia de la planificación de recursos humanos se mide a través de
un conjunto de indicadores clave (KPIs) que permiten detectar desviaciones
y tomar decisiones correctoras.

\subsubsection{KPIs utilizados por CaixaBank}

\begin{table}[H]
  \centering
  \small
  \caption{Indicadores de planificación de RRHH en CaixaBank}
  \begin{tabularx}{\textwidth}{p{0.35\textwidth} p{0.18\textwidth} X}
    \toprule
    \textbf{Indicador}                    & \textbf{¿Se mide?} & \textbf{Valor aproximado / fuente} \\
    \midrule
    Tasa de rotación (turnover)           & [Sí/No] & [\% — fuente: Memoria Anual / entrevista] \\
    Tasa de absentismo                    & [Sí/No] & [\%] \\
    Tiempo medio de cobertura (TTF)       & [Sí/No] & [días] \\
    Tasa de promoción interna             & [Sí/No] & [\%] \\
    Coste medio de contratación           & [Sí/No] & [\euro{}] \\
    Desviación plantilla real vs. plan    & [Sí/No] & [\%] \\
    Cobertura de sucesión (puestos clave) & [Sí/No] & [\%] \\
    ROI de formación                      & [Sí/No] & [\%] \\
    \bottomrule
  \end{tabularx}
\end{table}

[Análisis de los indicadores más relevantes obtenidos en la entrevista.
Para algunos datos públicos, consultad la Memoria Anual de CaixaBank.]

\subsubsection{Propuestas de mejora en los criterios de evaluación}

% TODO: ¿Qué KPIs deberían medir y no miden? Por ejemplo:
%   - NPS de empleados (Employee Net Promoter Score).
%   - Índice de diversidad e inclusión (si no lo tienen).
%   - Indicadores de bienestar (wellbeing index).
%   - Métricas de People Analytics predictivo.

\begin{itemize}
  \item [Propuesta 1 — justificación]
  \item [Propuesta 2 — justificación]
\end{itemize}

\subsection{Sistemas de Información de Recursos Humanos (SIRH)}

\subsubsection{¿Qué sistema SIRH utiliza CaixaBank?}

% TODO: CaixaBank, dada su escala, es probable que use uno de los grandes
% sistemas enterprise (SAP SuccessFactors, Workday, Oracle HCM).
% En 2021, tras la fusión, habrán tenido que integrar los sistemas de
% CaixaBank y Bankia. ¿Cuál es el sistema unificado?

[Descripción a completar con la información de la entrevista]

\begin{table}[H]
  \centering
  \small
  \caption{Módulos del SIRH de CaixaBank}
  \begin{tabularx}{\textwidth}{p{0.40\textwidth} p{0.12\textwidth} X}
    \toprule
    \textbf{Módulo / Funcionalidad}          & \textbf{¿Activo?} & \textbf{Observaciones} \\
    \midrule
    Administración de personal / nómina       & [Sí/No] & \\
    Gestión del tiempo y asistencia           & [Sí/No] & \\
    Reclutamiento y selección (ATS)           & [Sí/No] & \\
    Gestión del desempeño                     & [Sí/No] & \\
    Formación y desarrollo (LMS)              & [Sí/No] & \\
    Planificación de plantilla                & [Sí/No] & \\
    People Analytics e informes               & [Sí/No] & \\
    Gestión de la compensación                & [Sí/No] & \\
    Portal del empleado (autoservicio)        & [Sí/No] & \\
    \bottomrule
  \end{tabularx}
\end{table}

\subsubsection{¿Cómo mejorar el SIRH?}

% TODO: Una vez descripto el sistema actual, reflexionad sobre qué capacidades
% adicionales aportarían valor. Por ejemplo:
%   - Integrar People Analytics predictivo (predecir quién va a marcharse).
%   - Mejorar la experiencia del empleado en el portal de autoservicio.
%   - Integrar IA para el cribado de CV en selección (pero con control del sesgo).
%   - Mejorar la interoperabilidad entre el SIRH y otros sistemas del banco.

[Descripción a completar]

\begin{notaequipo}
  Preguntad en la entrevista:
  \begin{enumerate}
    \item ¿Qué sistema SIRH utilizan? ¿Ha cambiado tras la fusión con Bankia?
    \item ¿Qué módulos tienen activos? ¿Qué módulos les gustaría activar?
    \item ¿Qué KPIs de plantilla monitorizan de forma habitual?
    \item ¿Tienen un cuadro de mando (dashboard) de RRHH? ¿Quién lo consulta?
    \item ¿Están satisfechos con la información que les proporciona el SIRH?
          ¿Qué mejorarían?
  \end{enumerate}
\end{notaequipo}

\newpage

% ============================================================
%  TEMA 3 — APARTADO 5: PROCESOS SUSTRACTIVOS
%                        (Demanda < Oferta: exceso de plantilla)
%
%  RECORDATORIO AL EQUIPO:
%  Los procesos sustractivos son aquellos mediante los cuales la organización
%  reduce su plantilla cuando la oferta de RRHH supera la demanda.
%  Se clasifican en:
%
%  A) SUSPENSIÓN LABORAL TEMPORAL → ERTES
%     - Expediente de Regulación Temporal de Empleo.
%     - El contrato se SUSPENDE temporalmente (no se extingue).
%     - El trabajador cobra prestación del SEPE + posible complemento empresa.
%     - Causas: económicas, técnicas, organizativas, productivas, o FUERZA MAYOR.
%     - Muy usado durante la pandemia COVID-19.
%     - Ventaja: mantiene la plantilla formada cuando mejore la situación.
%     - Inconveniente: coste económico y de imagen; desmotivación.
%
%  B) RUPTURA LABORAL VOLUNTARIA
%     1. Dimisión (renuncia voluntaria del trabajador).
%        - El empleado abandona sin coste para la empresa (salvo sustitución).
%        - Puede ser signo de insatisfacción o de mejores oportunidades externas.
%     2. Jubilación (ordinaria, anticipada o parcial).
%        - Jubilación anticipada: la empresa incentiva la salida antes de la
%          edad legal, normalmente con complementos económicos.
%        - Es la forma más suave de reducir plantilla.
%
%  C) RUPTURA LABORAL INVOLUNTARIA
%     1. Despido individual (por causas objetivas o disciplinarias).
%     2. ERE — Expediente de Regulación de Empleo (despidos colectivos).
%        - Requiere periodo de consultas con representantes sindicales.
%        - Indemnizaciones + posibles medidas sociales.
%     3. Downsizing: reducción sistemática del tamaño organizativo.
%        - Puede acompañarse de reestructuración (rightsizing).
%     4. No renovación de contratos temporales.
%
%  D) RECOLOCACIÓN — OUTPLACEMENT
%     - Servicio de apoyo a los empleados que salen para que encuentren
%       nuevo empleo (orientación profesional, CV, entrevistas, networking).
%     - Puede ser gestionado internamente o a través de consultoras.
%     - Mejora la imagen de la empresa y reduce el trauma de la salida.
%
%  PARA ESTA SECCIÓN INDICAD:
%  - ¿Ha utilizado CaixaBank cada uno de estos mecanismos?
%  - Cuándo, bajo qué circunstancias y con qué resultados.
%  - Ventajas e inconvenientes de cada uno en el caso de CaixaBank.
%  - ¿Está la empresa satisfecha con cómo los ha gestionado?
% ============================================================

\section{Procesos sustractivos: gestión del exceso de plantilla}

Cuando la oferta de recursos humanos supera la demanda, CaixaBank debe
recurrir a mecanismos de reducción o reajuste de plantilla. A continuación
se analizan los principales procesos sustractivos aplicados o aplicables.

\subsection{Suspensión laboral temporal: los ERTE}

% TODO: ¿Ha recurrido CaixaBank a ERTE? Principalmente durante la pandemia
% COVID-19 (2020-2021). Si es así, ¿a cuántos empleados afectó? ¿Qué duración
% tuvo? ¿Qué complementos ofreció la empresa sobre la prestación del SEPE?
% ¿Fue eficaz para mantener el empleo y la operatividad del banco?

Un ERTE (Expediente de Regulación Temporal de Empleo) implica la suspensión
temporal del contrato de trabajo sin que se produzca su extinción. El
trabajador cobra una prestación por desempleo del SEPE, que la empresa puede
complementar.

\subsubsection{Aplicación en CaixaBank}

[Descripción a completar]

\textbf{Ventajas del ERTE para CaixaBank:}
\begin{itemize}
  \item [...]
\end{itemize}

\textbf{Inconvenientes / limitaciones:}
\begin{itemize}
  \item [...]
\end{itemize}

\subsection{Ruptura laboral voluntaria}

\subsubsection{Dimisión y rotación voluntaria}

% TODO: ¿Cuál es la tasa de dimisiones voluntarias en CaixaBank?
% ¿Realizan entrevistas de salida para conocer las causas?
% ¿Qué factores llevan a los empleados a marcharse voluntariamente
% (mejores salarios en la competencia, agotamiento, cambio de sector...)?

[Descripción a completar]

\subsubsection{Jubilación (ordinaria y anticipada)}

% TODO: CaixaBank tiene una plantilla relativamente envejecida en ciertos
% segmentos (herencia de Bankia y de la red de oficinas tradicionales).
% ¿Han implementado planes de jubilación anticipada?
% ¿En el ERE de 2021 se incluyeron medidas de jubilación anticipada?
% ¿Con qué condiciones (edad mínima, complemento sobre la pensión)?

[Descripción a completar]

\textbf{Ventajas de los planes de jubilación anticipada:}
\begin{itemize}
  \item Reducción no traumática de la plantilla sin despidos.
  \item Renovación generacional y entrada de perfiles más digitales.
  \item Menor conflicto sindical que los despidos colectivos.
\end{itemize}

\textbf{Inconvenientes:}
\begin{itemize}
  \item Pérdida de conocimiento experto y memoria institucional.
  \item Coste económico elevado para la empresa (complementos).
  \item Riesgo de que se jubilen empleados clave que no deberían marcharse.
\end{itemize}

\subsection{Ruptura laboral involuntaria: ERE y despidos}

\subsubsection{ERE 2021 — Fusión CaixaBank-Bankia}

% TODO: El ERE de 2021 fue el más grande de la historia bancaria española.
% Detalles clave a incluir:
%   - Número de afectados (aproximadamente 6.452 empleados según acuerdo sindical).
%   - Condiciones: indemnizaciones, bajas incentivadas, jubilaciones anticipadas.
%   - Sindicatos involucrados: UGT, CCOO, FINE.
%   - Reacción de los empleados y sindicatos.
%   - ¿Fue el ajuste suficiente o hay nuevos ajustes previstos?
%   - ¿Cómo afectó a la moral de la plantilla que se quedó?

En 2021, tras la fusión con Bankia, CaixaBank llevó a cabo el mayor
Expediente de Regulación de Empleo (ERE) de la banca española:

\begin{table}[H]
  \centering
  \small
  \caption{Datos del ERE CaixaBank 2021}
  \begin{tabularx}{\textwidth}{lX}
    \toprule
    \textbf{Parámetro}           & \textbf{Dato} \\
    \midrule
    Empleados afectados           & $\sim$6.452 (según acuerdo negociado) \\
    Cierre de oficinas            & $\sim$1.500 oficinas \\
    Tipo de salidas               & Bajas incentivadas, jubilaciones anticipadas, suspensiones \\
    Periodo de consultas          & [Fecha inicio – Fecha acuerdo] \\
    Sindicatos firmantes          & [CCOO, UGT, FINE — verificar] \\
    Medidas sociales de acompañamiento & [Outplacement, complementos, formación] \\
    \bottomrule
  \end{tabularx}
\end{table}

\textbf{Valoración del proceso:}

[Descripción a completar — ¿fue percibido como justo? ¿Qué impacto tuvo
en el clima laboral de la plantilla que se quedó? ¿Survivor syndrome?]

\subsubsection{Downsizing y reestructuración continua}

% TODO: Más allá del ERE de 2021, ¿hay un proceso de reducción continua
% (no renovación de contratos, digitalización que sustituye puestos, etc.)?
% La automatización de procesos bancarios y la banca digital reducen la
% necesidad de ciertos perfiles de back-office.

[Descripción a completar]

\subsection{Recolocación: el \textit{outplacement}}

\subsubsection{Servicios de \textit{outplacement} en el ERE de CaixaBank}

% TODO: ¿Ofreció CaixaBank servicios de outplacement a los empleados
% que salieron? ¿Lo gestionaron internamente o con una consultora externa
% (p.ej. Lee Hecht Harrison, Randstad RiseSmart, Adecco)?
% ¿Cuántos empleados se beneficiaron? ¿Con qué tasa de recolocación?

El outplacement es un conjunto de servicios de apoyo a los empleados
que abandonan la empresa (voluntaria o involuntariamente), orientados a
facilitar su reincorporación al mercado laboral:

\begin{itemize}
  \item Orientación y coaching profesional.
  \item Actualización y elaboración del currículum.
  \item Preparación para entrevistas de trabajo.
  \item Identificación de oportunidades en el mercado.
  \item Apoyo psicológico durante la transición.
\end{itemize}

[Descripción de cómo se aplicó en CaixaBank]

\textbf{Ventajas del outplacement para CaixaBank:}
\begin{itemize}
  \item Mejora la imagen como empleador responsable (\textit{employer branding}).
  \item Reduce el conflicto y los litigios laborales.
  \item Mitiga el impacto emocional en los empleados que salen, lo que
        también reduce el síndrome del superviviente en los que se quedan.
\end{itemize}

\begin{notaequipo}
  Preguntad en la entrevista:
  \begin{enumerate}
    \item ¿Han recurrido a ERTE en algún momento (COVID, fusión)?
    \item ¿Qué condiciones se ofrecieron en el ERE de 2021?
          ¿Hubo jubilaciones anticipadas? ¿A partir de qué edad?
    \item ¿Se ofreció outplacement a los empleados afectados por el ERE?
          ¿Quién lo gestionó? ¿Cuáles fueron los resultados?
    \item ¿Cómo afectó el ERE al clima laboral de los empleados que se
          quedaron? ¿Tomaron medidas para gestionar ese impacto?
    \item ¿Realizan entrevistas de salida cuando un empleado dimite?
          ¿Qué hacen con esa información?
  \end{enumerate}
\end{notaequipo}

\newpage

% ============================================================
%  TEMA 3 — APARTADO 6: ALTERNATIVAS A LA REDUCCIÓN DE PLANTILLA
%                        E INCORPORACIÓN A LA EMPRESA
%
%  RECORDATORIO AL EQUIPO:
%  Antes de recurrir a despidos o EREs, las empresas pueden adoptar
%  medidas alternativas para equilibrar oferta y demanda de RRHH:
%
%  ALTERNATIVAS A LA REDUCCIÓN (cuando hay superávit):
%    1. Reducción de la jornada laboral (trabajo a tiempo parcial).
%    2. Reparto del trabajo (job sharing / worksharing).
%    3. Reducción salarial temporal (sacrificio compartido).
%    4. Congelación de contrataciones y de salarios.
%    5. Formación y recualificación (reskilling): reconvertir empleados
%       en perfiles demandados (p.ej. de cajero a asesor digital).
%    6. Movilidad funcional: asignar empleados a otras áreas con demanda.
%    7. Movilidad geográfica: trasladar empleados a zonas con necesidades.
%    8. Excedencias voluntarias y licencias sin sueldo.
%    9. Reducción de horas extra.
%   10. Externalización de funciones no estratégicas (outsourcing) para
%       liberar capacidad interna hacia tareas de mayor valor.
%
%  ALTERNATIVAS A LA INCORPORACIÓN (cuando hay déficit):
%  En lugar de contratar externamente, la empresa puede:
%    1. Promoción interna (movilidad vertical).
%    2. Formación acelerada para cubrir el gap de habilidades.
%    3. Trabajo en equipo / colaboración cross-funcional.
%    4. Horas extra (si el déficit es temporal y limitado).
%    5. Trabajo temporal o contratos de proyecto.
%    6. Externalización (BPO, outsourcing, contratas).
%    7. Automatización (RPA, IA) para reducir la necesidad de personal.
%    8. Acuerdos de colaboración con otras empresas del sector.
%    9. Internalización de funciones previamente externalizadas.
%   10. Programas de recontratación de exempleados (boomerang employees).
%
%  PARA ESTA SECCIÓN INDICAD:
%  - ¿Qué alternativas ha utilizado o utiliza CaixaBank?
%  - ¿Han sido efectivas? Ventajas e inconvenientes concretos.
%  - ¿Qué alternativas serían interesantes para CaixaBank y no usa aún?
% ============================================================

\section{Alternativas a la reducción de plantilla e incorporación a la empresa}

\subsection{Alternativas a la reducción de plantilla}

Antes de recurrir a medidas de ruptura laboral, CaixaBank puede —y debe—
explorar alternativas menos traumáticas que permitan ajustar la plantilla
preservando el empleo y el clima organizacional.

\begin{table}[H]
  \centering
  \small
  \caption{Alternativas a la reducción de plantilla en CaixaBank}
  \begin{tabularx}{\textwidth}{p{0.35\textwidth} p{0.12\textwidth} X}
    \toprule
    \textbf{Alternativa}                        & \textbf{¿Se aplica?} & \textbf{Observaciones} \\
    \midrule
    Reducción de jornada / trabajo parcial       & [Sí/No] & \\
    Congelación de contrataciones                & [Sí/No] & \\
    Reskilling / reconversión de perfiles        & [Sí/No] & \\
    Movilidad funcional interna                  & [Sí/No] & \\
    Movilidad geográfica                         & [Sí/No] & \\
    Excedencias voluntarias incentivadas         & [Sí/No] & \\
    Outplacement previo al despido               & [Sí/No] & \\
    Reducción de horas extra / subcontratación   & [Sí/No] & \\
    \bottomrule
  \end{tabularx}
\end{table}

\subsubsection{Reskilling y recualificación profesional}

% TODO: Este es probablemente el aspecto más relevante para CaixaBank.
% La digitalización bancaria hace que muchos perfiles de red de oficinas
% queden obsoletos. ¿Tiene CaixaBank un programa de reskilling para
% reconvertir a estos empleados en perfiles digitales?
% Mencionar si existe un programa concreto (p.ej. "Future Skills" o similar).

[Descripción a completar]

\textbf{Ventajas:}
\begin{itemize}
  \item Preserva el empleo y el conocimiento del negocio.
  \item Mejor imagen como empresa responsable.
  \item Menor coste que despedir y contratar externamente.
\end{itemize}

\textbf{Inconvenientes:}
\begin{itemize}
  \item No todos los empleados tienen el potencial o la motivación para
        reconvertirse en perfiles muy técnicos.
  \item Requiere inversión elevada en formación.
  \item El proceso es lento; no resuelve necesidades inmediatas de
        perfiles digitales especializados.
\end{itemize}

\subsubsection{Movilidad interna (funcional y geográfica)}

% TODO: ¿Tiene CaixaBank un sistema formal de movilidad interna?
% ¿Los empleados pueden postularse a vacantes internas? ¿Hay un portal
% de oportunidades internas en el SIRH? ¿Con qué frecuencia se producen
% traslados geográficos voluntarios o forzosos?

[Descripción a completar]

\subsection{Alternativas a la incorporación de nueva plantilla}

Cuando existe déficit de personal, CaixaBank no siempre recurre de inmediato
a la contratación externa. Las alternativas disponibles incluyen:

\subsubsection{Promoción interna}

% TODO: ¿Cuál es la política de CaixaBank respecto a la promoción interna?
% ¿Se publican primero las vacantes internamente antes de salir al mercado?
% ¿Tienen planes de carrera definidos?

[Descripción a completar]

\subsubsection{Automatización y digitalización}

% TODO: La RPA (Automatización Robótica de Procesos) y la IA están
% transformando el sector bancario. CaixaBank ha invertido fuertemente en
% tecnología. ¿Hay procesos que antes requerían empleados y ahora están
% automatizados? ¿Esto ha reducido la necesidad de incorporaciones en
% ciertas áreas?

[Descripción a completar]

\subsubsection{Externalización (\textit{outsourcing})}

% TODO: ¿Qué funciones tiene CaixaBank externalizadas?
% ¿IT, call centers, servicios de back-office, seguridad, limpieza?
% ¿Han seguido la tendencia de insourcing (recuperar funciones) en algún área?

[Descripción a completar]

\subsubsection{Trabajo temporal y contratos de proyecto}

[Descripción a completar]

\begin{notaequipo}
  Preguntad en la entrevista:
  \begin{enumerate}
    \item ¿Tienen algún programa de reskilling o upskilling para reconvertir
          perfiles tradicionales? ¿En qué consiste? ¿Cuántos empleados
          se han beneficiado?
    \item ¿Existe un portal de vacantes internas para la movilidad interna?
    \item ¿Cómo deciden entre cubrir una vacante internamente o acudir
          al mercado externo?
    \item ¿Qué procesos han automatizado recientemente con RPA o IA?
          ¿Ha supuesto una reducción de la plantilla en esas áreas?
    \item ¿Qué funciones tienen externalizadas? ¿Han internalizado alguna
          función en los últimos años?
  \end{enumerate}
\end{notaequipo}

\newpage

% ============================================================
%  ANEXO: PREGUNTAS Y RESPUESTAS DE LA ENTREVISTA
% ============================================================
\input{secciones/12_anexo}

\end{document}
